%% LyX 2.4.4 created this file.  For more info, see https://www.lyx.org/.
%% Do not edit unless you really know what you are doing.
\documentclass[english]{article}
\usepackage[T1]{fontenc}
\usepackage[latin9]{inputenc}
\usepackage{enumitem}
\usepackage{amsmath}
\usepackage{amsthm}
\usepackage{cancel}
\usepackage{geometry}
\geometry{verbose,tmargin=1in,bmargin=1in,lmargin=1in,rmargin=1in}

\makeatletter

%%%%%%%%%%%%%%%%%%%%%%%%%%%%%% LyX specific LaTeX commands.
%% Because html converters don't know tabularnewline
\providecommand{\tabularnewline}{\\}

%%%%%%%%%%%%%%%%%%%%%%%%%%%%%% Textclass specific LaTeX commands.
\numberwithin{equation}{section}
\numberwithin{figure}{section}
\newlength{\lyxlabelwidth}      % auxiliary length 

%%%%%%%%%%%%%%%%%%%%%%%%%%%%%% User specified LaTeX commands.
\usepackage{hyperref}
\usepackage[usenames,dvipsnames,table]{xcolor} %adds additional color options
\hypersetup{colorlinks=true, urlcolor=black,linkcolor=black}

\usepackage{fancyhdr}
\usepackage{lastpage}
\pagestyle{fancy}
\usepackage{enumitem}
\usepackage{ifthen}
%sets math fonts to be more readable
\everymath=\expandafter{\the\everymath\displaystyle}
%\DeclareMathSizes{10}{10}{10}{10}
\usepackage{multicol}

\usepackage{tikz}
\usetikzlibrary{plotmarks}
\usepackage{pgfplots}
\pgfplotsset{compat=newest}

\usepackage{calc}
\usepackage[nomessages]{fp}

\usepackage{advdate}    % Advancing/saving dates
\usepackage{datetime}   % Dates formatting
\usepackage{datenumber} % Counters for dates

\usepackage{fancybox} %%ovalbox and fancy box settings

\usepackage[outline]{contour}  %allows text halos
\usepackage{colortbl}


\usepackage[super]{nth} %easy 1st, 2nd, etc

\usepackage{forest} %%for tree diagram

%%data for histogram%%
\begin{filecontents}{hist1_data.csv} 
dist 
3
1
0
4
1
3
2
2
0
2
0
2
2
1
4
3
1
1
3
4
2
1
3
0
1
0
2
5
1
2
3
0
0
1
2
3
1
2
0
2
\end{filecontents}

%%data for histogra: cal per cereal%%
\begin{filecontents}{hist_data_cereal.csv} 
dist 
3.7
3.6
3.6
4.0
4.0
3.6
3.8
3.7
4.1
3.9
3.9
3.3
3.5
3.6
3.9
4.0
\end{filecontents}
% Added by lyx2lyx
\setlength{\parskip}{\smallskipamount}
\setlength{\parindent}{0pt}

\makeatother

\usepackage{babel}
\begin{document}
\renewcommand{\author}{Dr.\,James Ashe}

\newcommand{\classNum}{107}

\newcommand{\classSec}{7}

\newcommand{\nameType}{}

\newcommand{\examType}{Final Exam Review}

\renewcommand{\date}{The Final is on \formatdate{5}{12}{2014} from 8am-10am}

%%First page's header%%%%%%%%%%%%%%%%%%%%%%
\fancypagestyle{firststyle} %%header settings for first pagr
{
	\fancyhead[L]{MTH \classNum{}(\classSec{}) \author{}\\
	\textbf{\examType{}} \date{}}
	\fancyhead[C]{\textbf{\nameType}}
	\setlength{\headsep}{14pt}
}
%%%%%%%%%%%%%%%%%%%%%%%%%%%%%%%%%%%%%%%%%%

%%Next page's header%%%%%%%%%%%%%%%%%%%%%%%
%\setlist{nolistsep}
\lhead{MTH \classNum{}(\classSec{})} 
\chead{\textbf{\nameType}} 
\cfoot{\thepage{}~of~\pageref{LastPage}} 
\setlength{\headsep}{5pt} 
%%%%%%%%%%%%%%%%%%%%%%%%%%%%%%%%%%%%%%%%%%%

%%Various settings; see the preamble for other settings%%%%%
%%\setenumerate[0]{leftmargin=12pt,itemindent=0pt,labelwidth=0pt}
%\setlist[2]{noitemsep}
%\setlist{nosep}
\setlist{itemsep=0pt,topsep=0pt}
\renewcommand{\labelenumi}{\textbf{\arabic{enumi}.}}
\renewcommand{\labelenumii}{\textbf{\alph{enumii}.}}
\thispagestyle{firststyle}
%%%%%%%%%%%%%%%%%%%%%%%%%%%%%%%%%%%%%%%%%%
\newcommand{\xmax}{10}
\newcommand{\xmin}{-10}
\newcommand{\xstep}{0} %must be xmin+n
\newcommand{\ymax}{10}
\newcommand{\ymin}{-10}
\newcommand{\ystep}{0} %must be ymin+n

\newcommand{\xspace}{\hspace*{40pt}}
\newcommand{\yspace}{\vspace*{.25cm}} %1.4cm
\newcommand{\rowColor}{gray!35}
\large\newcommand{\T}{}
\newcommand{\F}{}

\newcommand{\gauss}[2]{1/(#2*sqrt(2*pi))*exp(-((x-#1)^2)/(2*#2^2))}

\newcommand{\normalcdf}[4]{ %mean,std,z-left,z-right
\begin{tikzpicture}[scale=1.25]

\pgfmathsetmacro{\mn}{#1}
\pgfmathsetmacro{\sd}{#2}
\pgfmathsetmacro{\za}{#1+#2*1}
\pgfmathsetmacro{\zb}{#1+#2*2}
\pgfmathsetmacro{\zc}{#1+#2*3}

\pgfmathsetmacro{\zna}{#1-#2*1}
\pgfmathsetmacro{\znb}{#1-#2*2}
\pgfmathsetmacro{\znc}{#1-#2*3}

\pgfmathsetmacro{\zleft}{#3}
\pgfmathsetmacro{\zright}{#4}
\pgfmathsetmacro{\zsL}{(#3-#1)/#2} %z-score
\pgfmathsetmacro{\zsR}{(#4-#1)/#2} %z-score
\pgfmathsetmacro{\xCenter}{\zsL/2.5+\zsR/2.5} %center using z-score
\pgfmathsetmacro{\yCenter}{((1/(1*sqrt(2*pi))*exp(-((\xCenter-0)^2)/(2*1^2)))/2}
\pgfkeys{/pgf/fpu} %for large numbers
\pgfmathparse{100*(1/(1 + exp(-0.07056*((\zsR-0)/1)^3 - 1.5976*(\zsR-0)/1))-1/(1 + exp(-0.07056*((\zsL-0)/1)^3 - 1.5976*(\zsL-0)/1)))} %area from z-left to z-right
\edef\area{\pgfmathresult} 
\pgfkeys{/pgf/fpu=false}

\scriptsize
\begin{axis}[ 
	no markers, 
	domain=0:10, 
	%samples=100, 
	axis lines*=left, 
	%xlabel=Standard deviations, 
	%ylabel=Frequency, 
	height=2.5cm, %change/remove this scaling to get 1:1 pic
	width=5.5cm, 
	hide y axis,
	%ticks=none,
	xtick={-3,-2,-1,0,1,2,3},
	xticklabels={\pgfmathprintnumber[precision=0]{\znc},\pgfmathprintnumber[precision=2]{\znb},\pgfmathprintnumber[precision=2]{\zna},#1,\pgfmathprintnumber[precision=0]{\za},\pgfmathprintnumber[precision=2]{\zb},\pgfmathprintnumber[precision=2]{\zc}},
	%xticklabels={-25, -15, -5, 5, 15, 25, 35},
	%ytick=empty, 
	enlargelimits=false, 
	clip=false, 
	axis on top, 
	%grid = major
	]

	\addplot [fill=blue!35, draw=red,  domain=\zsL:\zsR] {\gauss{0}{1}} \closedcycle;
		\node (arrowEnd) at (axis cs: \xCenter, \yCenter) {};
	\addplot [draw, thick, smooth, domain=-3:3] {\gauss{0}{1}} \closedcycle;
		\node[anchor=south east] (arrowStart) at (current bounding box.east) {\pgfmathprintnumber[fixed,precision=1]{\area}\%};
	\draw[->,thick,black] (arrowStart) to (arrowEnd);
	
	%\addplot [fill=orange!20, draw=none, domain=2:3] {\gauss{0}{1}} \closedcycle; 
	%\addplot [fill=blue!20, draw=none, domain=-2:-1] {\gauss{0}{1}} \closedcycle; 
	%\addplot [fill=blue!20, draw=none, domain=1:2] {\gauss{0}{1}} \closedcycle; 
	%\only<1->{\node[above] at (axis cs: -0.5, 0.025) {$34.1\%$};} 
	%\only<1->{\node[above] at (axis cs: 0.5, 0.025) {$34.1\%$};} 
	%\only<1->{\node[above] at (axis cs: -1.5, 0.025) {$13.6\%$};}
	%\only<1->{\node[above] at (axis cs: 1.5, 0.025) {$13.6\%$};}
	%\only<1->{\node[above] at (axis cs: 2.5, 0.05) {$2.1\%$};}
	%\only<1->{\node[above] at (axis cs: -2.5, 0.05) {$2.1\%$};}
	%\only<1->{ \addplot[<->,black,very thick] coordinates {(-1,0.4) (1,0.4)};} 
	%\only<1->{\node[above] at (axis cs: 0, 0.4) {$68.2\%$};} 
	%\only<1->{\addplot[<->,black,very thick] coordinates {(-2,0.3) (2,0.3)};}
	%\only<1->{\node[above] at (axis cs: 0, 0.3) {$95.4\%$};}
	%\only<1->{\addplot[<->,black, very thick] coordinates {(-3,0.2) (3,0.2)};} 
	%\only<1->{\node[above] at (axis cs: 0, 0.2) {$99.7\%$};}
	\end{axis} 
\end{tikzpicture}
}

\newcommand{\normalpdf}[5]{ %mean,std,z-left,z-right
\begin{tikzpicture}[scale=1.25]

\pgfmathsetmacro{\mn}{#1}
\pgfmathsetmacro{\sd}{#2}
\pgfmathsetmacro{\za}{#1+#2*1}
\pgfmathsetmacro{\zb}{#1+#2*2}
\pgfmathsetmacro{\zc}{#1+#2*3}

\pgfmathsetmacro{\zna}{#1-#2*1}
\pgfmathsetmacro{\znb}{#1-#2*2}
\pgfmathsetmacro{\znc}{#1-#2*3}

\pgfmathsetmacro{\zleft}{#3}
\pgfmathsetmacro{\zright}{#4}
\pgfmathsetmacro{\zsL}{(#3-#1)/#2} %z-score
\pgfmathsetmacro{\zsR}{(#4-#1)/#2} %z-score
\pgfmathsetmacro{\zsV}{(#5-#1)/#2} %z-score
\pgfmathsetmacro{\xCenter}{\zsL/2.5+\zsR/2.5} %center using z-score
\pgfmathsetmacro{\yCenter}{((1/(1*sqrt(2*pi))*exp(-((\xCenter-0)^2)/(2*1^2)))/2}
\pgfkeys{/pgf/fpu} %for large numbers
\pgfmathparse{100*(1/(1 + exp(-0.07056*((\zsR-0)/1)^3 - 1.5976*(\zsR-0)/1))-1/(1 + exp(-0.07056*((\zsL-0)/1)^3 - 1.5976*(\zsL-0)/1)))} %area from z-left to z-right
\edef\area{\pgfmathresult} 
\pgfkeys{/pgf/fpu=false}

\scriptsize
\begin{axis}[ 
	no markers, 
	domain=0:10, 
	%samples=100, 
	axis lines*=left, 
	%xlabel=Standard deviations, 
	%ylabel=Frequency, 
	height=2.5cm, %change/remove this scaling to get 1:1 pic
	width=5.5cm, 
	hide y axis,
	%ticks=none,
	xtick={\zsV},
	xticklabels={\pgfmathprintnumber[precision=2]{#5}},
	%xticklabels={-25, -15, -5, 5, 15, 25, 35},
	%ytick=empty, 
	enlargelimits=false, 
	clip=false, 
	axis on top, 
	%grid = major
	]

	\addplot [fill=blue!35, draw=red,  domain=\zsL:\zsR] {\gauss{0}{1}} \closedcycle;
		\node (arrowEnd) at (axis cs: \xCenter, \yCenter) {};
	\addplot [draw, thick, smooth, domain=-3:3] {\gauss{0}{1}} \closedcycle;
		\node[anchor=south east] (arrowStart) at (current bounding box.east) {\pgfmathprintnumber[fixed,precision=1]{\area}\%};
	\draw[->,thick,black] (arrowStart) to (arrowEnd);
	
	%\addplot [fill=orange!20, draw=none, domain=2:3] {\gauss{0}{1}} \closedcycle; 
	%\addplot [fill=blue!20, draw=none, domain=-2:-1] {\gauss{0}{1}} \closedcycle; 
	%\addplot [fill=blue!20, draw=none, domain=1:2] {\gauss{0}{1}} \closedcycle; 
	%\only<1->{\node[above] at (axis cs: -0.5, 0.025) {$34.1\%$};} 
	%\only<1->{\node[above] at (axis cs: 0.5, 0.025) {$34.1\%$};} 
	%\only<1->{\node[above] at (axis cs: -1.5, 0.025) {$13.6\%$};}
	%\only<1->{\node[above] at (axis cs: 1.5, 0.025) {$13.6\%$};}
	%\only<1->{\node[above] at (axis cs: 2.5, 0.05) {$2.1\%$};}
	%\only<1->{\node[above] at (axis cs: -2.5, 0.05) {$2.1\%$};}
	%\only<1->{ \addplot[<->,black,very thick] coordinates {(-1,0.4) (1,0.4)};} 
	%\only<1->{\node[above] at (axis cs: 0, 0.4) {$68.2\%$};} 
	%\only<1->{\addplot[<->,black,very thick] coordinates {(-2,0.3) (2,0.3)};}
	%\only<1->{\node[above] at (axis cs: 0, 0.3) {$95.4\%$};}
	%\only<1->{\addplot[<->,black, very thick] coordinates {(-3,0.2) (3,0.2)};} 
	%\only<1->{\node[above] at (axis cs: 0, 0.2) {$99.7\%$};}
	\end{axis} 
\end{tikzpicture}
}\newcommand{\normalpdfB}[4]{ %mean,std,z-left,z-right
\begin{tikzpicture}[scale=1.25]

\pgfmathsetmacro{\mn}{#1}
\pgfmathsetmacro{\sd}{#2}
\pgfmathsetmacro{\za}{#1+#2*1}
\pgfmathsetmacro{\zb}{#1+#2*2}
\pgfmathsetmacro{\zc}{#1+#2*3}

\pgfmathsetmacro{\zna}{#1-#2*1}
\pgfmathsetmacro{\znb}{#1-#2*2}
\pgfmathsetmacro{\znc}{#1-#2*3}

\pgfmathsetmacro{\zleft}{#3}
\pgfmathsetmacro{\zright}{#4}
\pgfmathsetmacro{\zsL}{(#3-#1)/#2} %z-score
\pgfmathsetmacro{\zsR}{(#4-#1)/#2} %z-score
\pgfmathsetmacro{\xCenter}{\zsL/2.5+\zsR/2.5} %center using z-score
\pgfmathsetmacro{\yCenter}{((1/(1*sqrt(2*pi))*exp(-((\xCenter-0)^2)/(2*1^2)))/2}
\pgfkeys{/pgf/fpu} %for large numbers
\pgfmathparse{100*(1/(1 + exp(-0.07056*((\zsR-0)/1)^3 - 1.5976*(\zsR-0)/1))-1/(1 + exp(-0.07056*((\zsL-0)/1)^3 - 1.5976*(\zsL-0)/1)))} %area from z-left to z-right
\edef\area{\pgfmathresult} 
\pgfkeys{/pgf/fpu=false}

\scriptsize
\begin{axis}[ 
	no markers, 
	domain=0:10, 
	%samples=100, 
	axis lines*=left, 
	%xlabel=Standard deviations, 
	%ylabel=Frequency, 
	height=2.5cm, %change/remove this scaling to get 1:1 pic
	width=5.5cm, 
	hide y axis,
	%ticks=none,
	xtick={\zsL,\zsR},
	xticklabels={\pgfmathprintnumber[precision=2]{\zsL},\pgfmathprintnumber[precision=2]{\zsR}},
	%xticklabels={-25, -15, -5, 5, 15, 25, 35},
	%ytick=empty, 
	enlargelimits=false, 
	clip=false, 
	axis on top, 
	%grid = major
	]

	\addplot [fill=blue!35, draw=red,  domain=\zsL:\zsR] {\gauss{0}{1}} \closedcycle;
		\node (arrowEnd) at (axis cs: \xCenter, \yCenter) {};
	\addplot [draw, thick, smooth, domain=-3:3] {\gauss{0}{1}} \closedcycle;
		\node[anchor=south east] (arrowStart) at (current bounding box.east) {\pgfmathprintnumber[fixed,precision=1]{\area}\%};
	\draw[->,thick,black] (arrowStart) to (arrowEnd);
	
	%\addplot [fill=orange!20, draw=none, domain=2:3] {\gauss{0}{1}} \closedcycle; 
	%\addplot [fill=blue!20, draw=none, domain=-2:-1] {\gauss{0}{1}} \closedcycle; 
	%\addplot [fill=blue!20, draw=none, domain=1:2] {\gauss{0}{1}} \closedcycle; 
	%\only<1->{\node[above] at (axis cs: -0.5, 0.025) {$34.1\%$};} 
	%\only<1->{\node[above] at (axis cs: 0.5, 0.025) {$34.1\%$};} 
	%\only<1->{\node[above] at (axis cs: -1.5, 0.025) {$13.6\%$};}
	%\only<1->{\node[above] at (axis cs: 1.5, 0.025) {$13.6\%$};}
	%\only<1->{\node[above] at (axis cs: 2.5, 0.05) {$2.1\%$};}
	%\only<1->{\node[above] at (axis cs: -2.5, 0.05) {$2.1\%$};}
	%\only<1->{ \addplot[<->,black,very thick] coordinates {(-1,0.4) (1,0.4)};} 
	%\only<1->{\node[above] at (axis cs: 0, 0.4) {$68.2\%$};} 
	%\only<1->{\addplot[<->,black,very thick] coordinates {(-2,0.3) (2,0.3)};}
	%\only<1->{\node[above] at (axis cs: 0, 0.3) {$95.4\%$};}
	%\only<1->{\addplot[<->,black, very thick] coordinates {(-3,0.2) (3,0.2)};} 
	%\only<1->{\node[above] at (axis cs: 0, 0.2) {$99.7\%$};}
	\end{axis} 
\end{tikzpicture}
}\small
\begin{enumerate}[resume]
\item Using the symbolic representations\\
\begin{tabular}{rl}
$b$: & I ride a bicycle to school.\tabularnewline
$s$: & It is snowing.\tabularnewline
$r$: & It is raining. \tabularnewline
\end{tabular}\\
express the following statements in symbolic form.

\begin{enumerate}
\item It is raining and I rode my bicycle to school.

\begin{enumerate}
\item []$r\wedge b$\yspace\vfill
\end{enumerate}
\item If it is raining or snowing, then I will not ride my bicycle to school.

\begin{enumerate}
\item []$(r\vee s)\rightarrow\sim b$\yspace\vfill
\end{enumerate}
\end{enumerate}
\item Using the statement, ``If your order is over \$50, then you will
not have to pay shipping fees.'' and parts (a), (b), and (c) below,
complete the following:\emph{}\\
\emph{}

\begin{enumerate}
\item If $\mbox{you do not have to pay shipping fees}$ then $\mbox{your order is over \$50}$.
\item If $\mbox{your order is not over \$50}$ then $\mbox{you have to pay shipping fees}$.
\item If $\mbox{you have to pay shipping fees}$ then $\mbox{your order is not over \$50}$.
\item []Identify the converse of the statement. \underline{\hphantom{part a}}
\item []Identify the inverse of the statement. \underline{\hphantom{part a}}
\item []Identify the contrapositive of the statement. \underline{\hphantom{part a}}\vfill
\end{enumerate}
\item For parts \ref{enu:one}-\ref{enu:last}, identify the negation of
the provided statement.\newcommand{\negate}{students~}\begin{multicols}{2}

\begin{enumerate}
\item \label{enu:one}\renewcommand{\negate}{students~}Some \negate are
taking Math 107.

\begin{enumerate}
\item [(a)]All \negate are taking Math 107.
\item [(b)]No \negate are taking Math 107.
\item [(c)]Some \negate are not taking Math 107.
\item [(d)]Some \negate are taking Math 107.\vfill
\end{enumerate}
\item \renewcommand{\negate}{players~}All \negate need to attend practice.

\begin{enumerate}
\item [(a)]All \negate need to attend practice.
\item [(b)]No \negate need to attend practice.
\item [(c)]Some \negate do not need to attend practice.
\item [(d)]Some \negate need to attend practice.\vfill\columnbreak
\end{enumerate}
\item \renewcommand{\negate}{professors~}No \negate know how to read.

\begin{enumerate}
\item [(a)]All \negate know how to read.
\item [(b)]No \negate know how to read.
\item [(c)]Some \negate do not know how to read.
\item [(d)]Some \negate know how to read.\vfill
\end{enumerate}
\item \label{enu:last}\renewcommand{\negate}{students~}Some \negate are
going to lunch.

\begin{enumerate}
\item [(a)]All \negate are going to lunch.
\item [(b)]No \negate are going to lunch.
\item [(c)]Some \negate are not going to lunch.
\item [(d)]Some \negate are going to lunch.\vfill
\end{enumerate}
\end{enumerate}
\end{multicols}\vfill
\item Use a truth table to determine if the following arguments are valid.

\begin{enumerate}
\item \vspace{-25pt}%
\begin{tabular}{rll}
 &  & \tabularnewline
 &  & \tabularnewline
$b$: & You buy the bike. & 1. You buy the bike or the tent.\tabularnewline
$t$: & You buy the tent. & 2. You buy the bike.\tabularnewline
\cline{3-3}
 &  & Therefore, you do not buy the tent.\tabularnewline
\end{tabular}\\
\begin{tabular}{c|c|c|c|c|c|c}
\multicolumn{1}{c}{} & \multicolumn{1}{c}{} & \multicolumn{1}{c}{\xspace} & \multicolumn{1}{c}{\xspace} & \multicolumn{1}{c}{\xspace} & \multicolumn{1}{c}{\xspace} & \xspace\tabularnewline
\multicolumn{1}{c}{} & \multicolumn{1}{c}{} & \multicolumn{1}{c}{\ovalbox{$1$}} & \multicolumn{1}{c}{\ovalbox{$2$}} & \multicolumn{1}{c}{} & \multicolumn{1}{c}{\ovalbox{$c$}} & \ovalbox{Argument}\tabularnewline
$b$ & $t$ & $b\vee t$ & $b$ & $1 \wedge 2$ & $\sim t$ & $\left(1\wedge2\right)\rightarrow c$\tabularnewline
\hline 
\hline 
\rowcolor{\rowColor}T & T & \T & \T & \T & \F & \F\tabularnewline
\hline 
T & F & \T & \T & \T & \T & \T\tabularnewline
\hline 
\rowcolor{\rowColor}F & T & \T & \F & \F & \F & \T\tabularnewline
\hline 
F & F & \F & \F & \F & \T & \T\tabularnewline
\end{tabular}\vfill
\item \vspace{-25pt} %
\begin{tabular}{rll}
 &  & \tabularnewline
 &  & \tabularnewline
$r$: & You register early. & 1. If you register early, then you will receive a discount.\tabularnewline
$d$: & You receive a discount. & 2. You register early.\tabularnewline
\cline{3-3}
 &  & Therefore, you will receive a discount.\tabularnewline
\end{tabular}\\
\begin{tabular}{c|c|c|c|c|c|c}
\multicolumn{1}{c}{} & \multicolumn{1}{c}{} & \multicolumn{1}{c}{\xspace} & \multicolumn{1}{c}{\xspace} & \multicolumn{1}{c}{\xspace} & \multicolumn{1}{c}{\xspace} & \xspace\tabularnewline
\multicolumn{1}{c}{} & \multicolumn{1}{c}{} & \multicolumn{1}{c}{\ovalbox{$1$}} & \multicolumn{1}{c}{\ovalbox{$2$}} & \multicolumn{1}{c}{} & \multicolumn{1}{c}{\ovalbox{$c$}} & \ovalbox{Argument}\tabularnewline
$r$ & $d$ & $r\rightarrow d$ & $r$ & $1 \wedge 2$ & $r$ & $\left(1\wedge2\right)\rightarrow c$\tabularnewline
\hline 
\hline 
\rowcolor{\rowColor}T & T & \T & \T & \T & \T & \T\tabularnewline
\hline 
T & F & \F & \T & \F & \T & \T\tabularnewline
\hline 
\rowcolor{\rowColor}F & T & \T & \F & \F & \F & \T\tabularnewline
\hline 
F & F & \T & \F & \F & \F & \T\tabularnewline
\end{tabular}\vfill\newpage
\end{enumerate}
\item Let $U=\{1,2,3,4,5,6,7,8\}$ be the universal set, and $A=\{1,2,3,4\}$
and $B=\{4,6,8\}$ be subsets of $U$. Find the indicated set, and
draw a Venn diagram corresponding to the indicated set.\newcommand{\vennScale}{0.5}\begin{multicols}{2}

\begin{enumerate}
\item $A\cup B$\vfill
\item $A\cap B$\vfill
\item $A'$\vfill
\item $B'$\vfill
\item $A\cap B'$\vfill
\item $A'\cup B$\vfill
\item $(A\cap B)'$\vfill
\item $(A\cup B)'$\vfill
\end{enumerate}
\end{multicols}
\item A college has $1105$ students. This semester $320$ students are
taking a math course, $405$ of the students are taking a history
class, and $250$ of the students taking a math class are not enrolled
in a history class.

\begin{enumerate}
\item Construct a Venn diagram, and write the appropriate numbers in the
four regions
\item How many of the students are taking math and history?
\item How many students are taking history but not math?
\item How many students are taking math or history?
\item How many students are taking neither math nor history?\vfill
\end{enumerate}
\item \emph{Survivor} and \emph{CSI} belong to the Thursday night prime
time lineup. A group of $600$ college students were surveyed. $350$
students watch \emph{Survivor} and $280$ watch \emph{CSI}. $180$
students watch both shows.

\begin{enumerate}
\item Construct a Venn diagram, and write the appropriate numbers in the
eight regions.
\item How many students watch only \emph{Survivor}?
\item How many students watch \emph{CSI }or \emph{Survivor}?
\item How many students watch only \emph{CSI}?\vfill
\end{enumerate}
\item Find the value of each of the following:\begin{multicols}{2}

\begin{enumerate}
\item $4!$\yspace\vfill
\item $7!$\yspace\vfill
\item $0!$\yspace\vfill
\item $\frac{10!}{7!}$\yspace\vfill
\item $\frac{27!}{4!25!}$\yspace\vfill
\item $\frac{1000!}{998!}$\yspace\vfill
\item $_{7}P_{2}$\yspace\vfill
\item $_{6}C_{4}$\yspace\vfill
\end{enumerate}
\end{multicols}
\item A certain model of pickup truck is available in five exterior colors,
three interior colors, and three interior styles. In addition, the
transmission can be either manual or automatic, and the truck can
have either two-wheel or four-wheel drive. How many different versions
of the pickup truck can be ordered?\vfill
\item A certain fraternity must elect a president, vice president, and treasurer
from its 30 members. How many ways can this be done?\vfill
\item A three-person committee must be chosen from a group of ten men and
twelve women. How many committees are possible if it must consist
of the following?

\begin{enumerate}
\item Three men.
\item Three women.
\item One woman and two men.
\item One man and two women.
\item One man and two women, and one of these is selected as a spokesperson
of the committee.\vfill\newpage
\end{enumerate}
\item A card is dealt from a well-shuffled, standard deck of $52$ cards.
Find the probability of being dealt each of the following.\begin{multicols}{2}

\begin{enumerate}
\item A red card\vfill
\item A queen\vfill
\item A club\vfill
\item The queen of clubs\vfill
\item A queen or a club\vfill
\item not a queen\vfill
\end{enumerate}
\end{multicols}\vfill
\item Three coins are tossed. Find each of the following:\begin{multicols}{2}

\begin{enumerate}
\item the sample space\vfill
\item the event, $E_{1}$, that exactly two are tails.\vfill
\item the event, $E_{2}$, that at least two are tails.\vfill
\item the probability of $E_{1}$.
\item the probability $E_{2}$.
\end{enumerate}
\end{multicols}\vfill
\item A drawer has $7$ socks. $4$ socks are black and $3$ are white socks.
John randomly pulls out $4$ socks. Find the probability of the following:\begin{multicols}{2}

\begin{enumerate}
\item all $4$ socks are black.\vfill
\item exactly $2$ are white.\vfill

\columnbreak
\item at least $3$ are white.\vfill\vfill
\end{enumerate}
\end{multicols}\vfill
\item Of all the flashlights is a large shipment, $15\%$ have a defective
bulb, $10\%$ have a defective battery, and $5\%$ have both defects.
If you buy one of these flashlights, find the probability of the following:\begin{multicols}{2}

\begin{enumerate}
\item a defective bulb or a defective battery\vfill
\item a good bulb or a good battery\vfill
\item a defective bulb and a defective battery. \vfill
\item just a defective bulb\vfill

\begin{enumerate}
\item []\columnbreak
\end{enumerate}
\end{enumerate}
\end{multicols}\vfill
\item Use the provided table to answer the following questions.

\begin{tabular}{|c|c|c|c|}
\hline 
\textbf{Age} & \textbf{Obama} & \textbf{McCain} & \textbf{Other/NA}\tabularnewline
\hline 
18-29 & 11.9\% & 5.8\% & 0.4\%\tabularnewline
\hline 
30-44 & 15.1\% & 13.3\% & 0.6\%\tabularnewline
\hline 
45-64 & 18.5\% & 18.1\% & 0.4\%\tabularnewline
\hline 
65+ & 7.2\% & 8.5\% & 0.3\%\tabularnewline
\hline 
\end{tabular}
\begin{enumerate}
\item What is the probability that a voter under $45$ voted for McCain.
\vfill
\item What is the probability that a person who voted for McCain is $45$
or older.\vfill
\end{enumerate}
\vfill
\item Gregor's Garden Corner buys $40\%$ of its plants from the Green Growery
and the rest from Herb's Herbs. Of the plants from the Green Growery,
$20\%$ must be returned, and $10\%$ of those from Herb's Herbs must
be returned. \begin{multicols}{2}

\begin{enumerate}
\item draw a tree diagram describing the probability\vfill
\item find the probability that a plant must be returned and it is from
Herb's Herbs.\vfill
\item find the probability that a plant from the Green Growery must be returned.\vfill
\item find the probability that a plant must be returned. \vfill
\end{enumerate}
\end{multicols}\vfill\newpage
\item Sixteen randomly selected families were surveyed to determine the
number of children in each family. The following results were found:

\hfill%
\begin{tabular}{cccccccc}
3.7  & 3.6 & 4.0 & 3.8 & 4.1 & 3.9 & 3.5 & 3.9\tabularnewline
3.6 & 4.0 & 3.6 & 3.7 & 3.9 & 3.3 & 3.6 & 4.0\tabularnewline
\end{tabular}\hfill\hspace{0cm}\begin{multicols}{2}
\begin{enumerate}
\item Find the mean.\vfill
\item Find the median.\vfill
\item Find the standard deviation.\vfill
\item Construct a histogram using 5 subintervals, beginning at 3.2, and
ending at 4.2.\columnbreak
\item []\vspace{4cm}
\end{enumerate}
\end{multicols}\vfill
\item Find the following probabilities, 

\begin{enumerate}
\item $p(z<1.35)$\vfill
\item $p(0<z<2.02)$\vfill
\item $p(-.09<z<1.06)$\vfill
\item $p(-1.1<z<-.08)$\vfill
\end{enumerate}
\item A population is normally distributed with a mean 78 and standard deviation
7.5. Find the following probabilities. 

\begin{enumerate}
\item $p(x<70)$\vfill
\item $p(70<x<85)$\vfill
\item $p(80<x<95.5)$\vfill
\item $p(95.5<x<108)$\vfill
\end{enumerate}
\item Find $c$ such that each of the following is true.

\begin{enumerate}
\item $p(z<c)=0.2776$\vfill
\item $p(z>c)=.0351$\vfill
\item Let $\mu=100$ and $\sigma=15$ find $c$ such that $p(x<c)=27.76\%$\vfill
\end{enumerate}
\item The time it takes latex paint to dry is normally distributed with
$\mu=3.5$ hours and $\sigma=.75$ hours. Find the probability that
drying time will be as follows.

\begin{enumerate}
\item less than $2.25$ hours.\vfill
\end{enumerate}
\item The average score on the ACT is $18$ with a standard deviation of
$6$. If a college only wants to accept students with an ACT score
at least in the top 25\%. What is the minimum score they should accept?
What about the top 5\%?\vfill
\end{enumerate}

\end{document}
